%% Determines the type of document (including standard settings for layout).
\documentclass[letterpaper]{article}



%% Package to control the size of a page: the standard margins used by LaTeX are
%% very wide!
\usepackage[margin=1in]{geometry}


%% Packages add functionality to the document. The AMS packages are standard
%% packages to support various mathematical notations. AMS stands for American
%% Mathematical Society, the organization that maintains these packages.
\usepackage{amsmath,amsthm,amssymb}

%% Theorems-like environments using functionality provided by amsthm.
\theoremstyle{plain}
\newtheorem{theorem}{Theorem}[section]
    %% [section] at the end specifies that theorems should be numbered
    %% per-section: Section x starts with theorem-like x.1 and so on, ...
\newtheorem{proposition}[theorem]{Proposition}
    %% [theorem] in the middle specifies: use the same counter as the theorem
    %% environment: here we number all theorem-like environments consecutively.
\newtheorem{corollary}[theorem]{Corollary}
\newtheorem{lemma}[theorem]{Lemma}
\theoremstyle{definition}
\newtheorem{definition}[theorem]{Definition}
\theoremstyle{remark}
\newtheorem{example}[theorem]{Example}
\newtheorem{remark}[theorem]{Remark}


%% Support for nicely formatted tables.
\usepackage{booktabs}


%% Support for colors & colors in tables.
\usepackage[table]{xcolor}

% Seven colors safe for use color blindness.
% Colors taken from doi:10.1038/nmeth.1618.
\definecolor{cbOrange}{RGB}{230,159,0}
\definecolor{cbSkyBlue}{RGB}{86,180,233}
\definecolor{cbBluishGreen}{RGB}{0,158,115}
\definecolor{cbYellow}{RGB}{240,228,66}
\definecolor{cbBlue}{RGB}{0,114,178}
\definecolor{cbVermillion}{RGB}{213,94,0}
\definecolor{cbReddischPurple}{RGB}{204,121,167}


%% Notation used in this document.
\newcommand{\n}{\mathbf{n}} %% Num. Replicas.
\newcommand{\f}{\mathbf{f}} %% Num. Faulty Replicas.

%% Misc. Math notation.
\newcommand{\BigO}{\mathcal{O}}
\newcommand{\abs}[1]{\lvert#1\rvert}
\newcommand{\AName}[1]{\textsc{#1}}
\newcommand{\Var}[1]{\texttt{#1}}


%% Algorithms.
\usepackage{algorithm}
\usepackage[noend]{algorithmic}
\newcommand{\GETS}{:=}


%% Formatting SI-units.
\usepackage{siunitx}
\sisetup{per-mode=symbol}


%% TikZ: for creating figures.
\usepackage{tikz}

%% Configuration for figures: Nicer arrows.
\usetikzlibrary{arrows.meta}
\tikzset{>=Stealth}


%% pgfplots: drawing plots using TikZ.
\usepackage{pgfplots}
%% Configuration for plots: Use color-blind friendly colors.
\pgfplotscreateplotcyclelist{cbSafeList}{
    very thick,solid,cbOrange,every mark/.append style={solid},mark=*\\
    very thick,solid,cbSkyBlue,every mark/.append style={solid},mark=*\\
    very thick,solid,cbBluishGreen,every mark/.append style={solid},mark=*\\
    very thick,solid,cbYellow,every mark/.append style={solid},mark=*\\
    very thick,solid,cbBlue,every mark/.append style={solid},mark=*\\
    very thick,solid,cbVermillion,every mark/.append style={solid},mark=*\\
    very thick,solid,cbReddischPurple,every mark/.append style={solid},mark=*\\
    very thick,solid,black,every mark/.append style={solid},mark=*\\
}
\pgfplotsset{
    legend style={font=\small},
    compat=1.16,
    width=260pt,
    height=140pt,
    legend cell align=left,
    xlabel near ticks,
    ylabel near ticks,
    every axis/.append style={
        cycle list name=cbSafeList,
        ymin=0,
        enlargelimits=0.05,
        mark size=1pt,
        ylabel style={align=center},
        xlabel style={align=center},
        title style={align=center}
    }
}

%% PgfplotsTable: loading data files to use with pgfplots.
\usepackage{pgfplotstable}


%% Support for hyperlinks and urls. The setting ``colorlinks'' sets how links
%% are shown in the document (with a color, without underline). We put hyperref
%% last---it has a tendency to break other packages when loaded before them.
\usepackage[colorlinks]{hyperref}
\usepackage{graphicx}
\usepackage{pgffor}
\usepackage{caption}
\usepackage{tabularx}
\usepackage{enumitem}
\usepackage{fancyhdr}
\usepackage{comment}


\begin{document}

%% First the meta-data of the doucment: title,

\begin{titlepage}
    \centering
    \vspace*{3cm}
    
    {\Huge CS3MI3 Assignment 1}\\[2cm]
    
    {\large Luca Mawyin}\\[0.5cm]
    
    {\large Dr.\ JP Rooney}\\[0.5cm]

    {\large \today}\\[2cm]

    {\Large McMaster University}

    \vfill
    \begin{flushleft}

        {\large \textbf{Student Number: }400531739}\\[0.5cm]
        {\large \textbf{MacID:} mawyinl}
    \end{flushleft}

\end{titlepage}

\renewcommand{\thesection}{Question \arabic{section}}

\section{}

\textbf{Let} $t,t_0,t_1,t_2\in\text{arith}$, then:

\begin{align*}
    \text{esses}(0)&=0\\
    \text{esses}(\text{true})&=0\\
    \text{esses}(\text{false})&=1\\
    \text{esses}(\text{pred}~t)&=\text{esses}(t)\\
    \text{esses}(\text{succ}~t)&=1+\text{esses}(t)\\
    \text{esses}(\text{iszero}~t)&=1+\text{esses}(t)\\
    \text{esses}(\text{if $t_0$ then $t_1$ else $t_2$})&=1+\text{esses}(t_0)+\text{esses}(t_1)+\text{esses}(t_2)
\end{align*}

\section{}

Firstly, we will define $\text{size}(t),~t,t_0,t_1,t_2\in\text{arith}$:

\begin{align*}
    \text{size}(0)&=1\\
    \text{size}(\text{false})&=1\\
    \text{size}(\text{true})&=1\\
    \text{size}(\text{succ}~t)&=1+\text{size}(t)\\
    \text{size}(\text{pred}~t)&=1+\text{size}(t)\\
    \text{size}(\text{iszero}~t)&=1+\text{size}(t)\\
    \text{size}(\text{if $t_0$ then $t_1$ else $t_2$})&=1+\text{size}(t_0)+\text{size}(t_1)+\text{size}(t_2)
\end{align*}
Next, we will prove that the property holds from our base cases:\\[1em]
Case $t=0$
\begin{align*}
    \text{esses}(0)&\leq\text{size}(0)\\
    0&\leq1
\end{align*}
Case $t=\text{true}$
\begin{align*}
    \text{esses}(\text{true})&\leq\text{size}(\text{true})\\
    0&\leq1
\end{align*}
Case $t=\text{false}$
\begin{align*}
    \text{esses}(\text{false})&\leq\text{size}(\text{false})\\
    1&\leq1
\end{align*}

Therefore, the property holds for all bases cases. Now we will prove that the property holds for all inductive cases $t\in\text{arith}$:\\[1em]
Case $t=\text{succ}~t_0$\\[1em]
Assuming:
\begin{align*}
    \text{esses}(t_0)&\leq\text{size}(t_0)
\end{align*}
then,
\begin{align*}
    \text{esses}(\text{succ}~t_0)&\leq\text{size}(\text{succ}~t_0)\\
    1+\text{esses}(t_0)&\leq1+\text{size}(t_0)\\
    \text{esses}(t_0)&\leq\text{size}(t_0)
\end{align*}
Using the \textbf{induction hypothesis}, our property holds.\\[1em]
Case $t=\text{pred}~t_0$\\[1em]
Assuming:
\begin{align*}
    \text{esses}(t_0)&\leq\text{size}(t_0)
\end{align*}
then,
\begin{align*}
    \text{esses}(\text{pred}~t_0)&\leq\text{size}(\text{pred}~t_0)\\
    \text{esses}(t_0)&\leq1+\text{size}(t_0)\\
\end{align*}
As $\text{esses}(t_0)\leq\text{size}(t_0)$ by our assumption, $\text{esses}(t_0)\leq1+\text{size}(t_0)$ will hold, trivially.\\[1em]
Case $t=\text{iszero}~t_0$\\[1em]
Assuming:
\begin{align*}
    \text{esses}(t_0)&\leq\text{size}(t_0)
\end{align*}
then,
\begin{align*}
    \text{esses}(\text{iszero}~t_0)&\leq\text{size}(\text{iszero}~t_0)\\
    1+\text{esses}(t_0)&\leq1+\text{size}(t_0)\\
    \text{esses}(t_0)&\leq\text{size}(t_0)
\end{align*}
Using the \textbf{induction hypothesis}, our property holds.\\[1em]
Case $t=\text{if}~t_0~\text{then}~t_1~\text{else}~t_2$\\[1em]
Assuming:
\begin{align*}
    \text{esses}(t_0)&\leq\text{size}(t_0),\\
    \text{esses}(t_1)&\leq\text{size}(t_1),\\
    \text{esses}(t_2)&\leq\text{size}(t_2)
\end{align*}
then,
\begin{align*}
    \text{esses}(\text{if}~t_0~\text{then}~t_1~\text{else}~t_2)&\leq\text{size}(\text{if}~t_0~\text{then}~t_1~\text{else}~t_2)\\
    1+\text{esses}(t_0)+\text{esses}(t_1)+\text{esses}(t_2)&\leq1+\text{size}(t_0)+\text{size}(t_1)+\text{size}(t_2)\\
    \text{esses}(t_0)+\text{esses}(t_1)+\text{esses}(t_2)&\leq\text{size}(t_0)+\text{size}(t_1)+\text{size}(t_2)\\
\end{align*}
Using the \textbf{induction hypothesis}, our property holds.

\section{}
\begin{verbatim}
esses :: ArithTerm -> Int
esses TmTrue = 0
esses TmFalse = 1
esses TmZero = 0
esses (TmSucc t) = 1 + esses t
esses (TmPred t) = esses t
esses (TmIsZero t) = 1 + esses t
esses (TmIfThenElse t0 t1 t2) = 1 + esses t0 + esses t1 + esses t2
\end{verbatim}

\end{document}